\documentclass[11pt,a4paper]{article}
\usepackage[utf8]{inputenc}
\usepackage[T1]{fontenc}
\usepackage[spanish]{babel}
\usepackage{amsmath,amssymb,mathtools}
\usepackage{geometry}
\usepackage{enumitem}
\usepackage{lmodern}
\usepackage{microtype}
\geometry{margin=2.4cm}
\setlength{\parindent}{0pt}
\setlength{\parskip}{0.45em}
\newcommand{\R}{\mathbb{R}}
\newcommand{\tr}{\operatorname{tr}}
\newcommand{\diag}{\operatorname{diag}}
\begin{document}

\begin{center}
{\LARGE\bfseries Álgebra Lineal}\\[0.4em]
{\Large Tema 7: Cónicas y Cuádricas}
\end{center}

\tableofcontents
\newpage

\section{Cónicas}

Una cónica es el lugar geométrico de los puntos del plano que verifican una ecuación de segundo grado en dos variables.

\subsection{Elipse}

Una elipse es el lugar geométrico de los puntos cuya suma de distancias a dos focos es constante:

\[
d(P,F_1)+d(P,F_2)=2a.
\]

Tomando los focos:

\[
F_1=(c,0),
\qquad
F_2=(-c,0),
\]

y definiendo:

\[
b^2=a^2-c^2,
\]

se obtiene la ecuación reducida:

\[
\boxed{
\frac{x^2}{a^2}
+
\frac{y^2}{b^2}
=
1.
}
\]

\subsection{Hipérbola}

Una hipérbola es el lugar geométrico de los puntos cuya diferencia, en valor absoluto, de distancias a dos focos es constante:

\[
|d(P,F_1)-d(P,F_2)|=2a.
\]

Tomando:

\[
F_1=(c,0),
\qquad
F_2=(-c,0),
\]

y:

\[
b^2=c^2-a^2,
\]

se obtiene:

\[
\boxed{
\frac{x^2}{a^2}
-
\frac{y^2}{b^2}
=
1.
}
\]

\subsection{Parábola}

Una parábola es el lugar geométrico de los puntos que equidistan de un foco y una recta directriz.

Si:

\[
F=\left(\frac p2,0\right)
\]

y la directriz es:

\[
x=-\frac p2,
\]

entonces:

\[
d(P,F)=d(P,r).
\]

La ecuación reducida es:

\[
\boxed{
y^2=2px.
}
\]

\subsection{Secciones cónicas}

La elipse, la hipérbola y la parábola aparecen al cortar un cono mediante distintos planos.

También aparecen cónicas degeneradas:

\begin{itemize}
\item un punto;
\item dos rectas secantes;
\item una recta doble;
\item dos rectas paralelas;
\item dos rectas paralelas imaginarias.
\end{itemize}

\section{Ecuación general de una cónica}

La ecuación general es:

\[
\boxed{
a_{11}x^2+a_{22}y^2+2a_{12}xy+2a_1x+2a_2y+a_0=0.
}
\]

Definimos:

\[
X=
\begin{pmatrix}
x\\y
\end{pmatrix},
\qquad
A=
\begin{pmatrix}
a_{11}&a_{12}\\
a_{12}&a_{22}
\end{pmatrix},
\qquad
B=
\begin{pmatrix}
2a_1&2a_2
\end{pmatrix}.
\]

Entonces:

\[
\boxed{
X^tAX+BX+a_0=0.
}
\]

En coordenadas homogéneas:

\[
\widetilde X=
\begin{pmatrix}
1\\x\\y
\end{pmatrix},
\qquad
\widetilde A=
\begin{pmatrix}
a_0&a_1&a_2\\
a_1&a_{11}&a_{12}\\
a_2&a_{12}&a_{22}
\end{pmatrix},
\]

y:

\[
\boxed{
\widetilde X^t\widetilde A\widetilde X=0.
}
\]

\section{Ecuación reducida de una cónica}

La ecuación:

\[
X^tAX+BX+a_0=0
\]

está reducida cuando:

\begin{itemize}
\item $A$ es diagonal;
\item si $0$ no es autovalor de $A$, los términos lineales han desaparecido;
\item si $0$ es autovalor de $A$, queda como máximo un término lineal.
\end{itemize}

La reducción se realiza mediante un movimiento rígido: una rotación seguida de una traslación.

\subsection{Paso 1: rotación}

Como $A$ es real y simétrica, existe una matriz ortogonal $P$ tal que:

\[
\boxed{
D=P^tAP.
}
\]

Con el cambio:

\[
X=PX',
\]

se obtiene:

\[
\boxed{
X'^tDX'+BPX'+a_0=0.
}
\]

La matriz $D$ contiene los autovalores de $A$.

\subsection{Paso 2: traslación}

Los términos lineales se eliminan completando cuadrados.

Por ejemplo:

\[
x'^2+2bx'
=
(x'+b)^2-b^2.
\]

Se introduce una traslación:

\[
X''=C+X'.
\]

La composición de rotación y traslación proporciona el cambio de referencia que lleva la ecuación general a la reducida.

\section{Clasificación de cónicas por la ecuación reducida}

Si:

\[
\lambda_1x^2+\lambda_2y^2-c=0,
\]

entonces:

\subsection{Tipo elíptico}

Si:

\[
\lambda_1\lambda_2>0,
\]

se tiene:

\[
\boxed{
\begin{cases}
c=0 & \text{punto},\\
\lambda_1c>0 & \text{elipse real},\\
\lambda_1c<0 & \text{elipse imaginaria}.
\end{cases}
}
\]

\subsection{Tipo hiperbólico}

Si:

\[
\lambda_1\lambda_2<0,
\]

entonces:

\[
\boxed{
\begin{cases}
c\neq0 & \text{hipérbola},\\
c=0 & \text{dos rectas secantes}.
\end{cases}
}
\]

\subsection{Tipo parabólico}

Si:

\[
\lambda_1\lambda_2=0,
\]

y, por ejemplo, $\lambda_1=0$, tras diagonalizar se obtiene:

\[
\lambda_2y'^2+2b_1x'+2b_2y'+a_0=0.
\]

Si:

\[
b_1\neq0,
\]

es una parábola.

Si:

\[
b_1=0,
\]

se llega a:

\[
\lambda_2y''^2=c,
\]

que puede representar:

\begin{itemize}
\item dos rectas coincidentes;
\item dos rectas paralelas reales;
\item dos rectas paralelas imaginarias.
\end{itemize}

\section{Invariantes métricos de una cónica}

Se definen:

\[
\boxed{
I_3=\det\widetilde A
}
\]

\[
\boxed{
I_2=\det A
}
\]

\[
\boxed{
I_1=\tr A.
}
\]

Estos valores no cambian bajo movimientos rígidos.

La clasificación indicada en la presentación es:

\[
\boxed{
I_3\neq0,\ I_2>0
\Rightarrow
\text{elipse}
}
\]

y dentro de este caso:

\[
I_1I_3<0
\Rightarrow
\text{elipse real},
\]

\[
I_1I_3>0
\Rightarrow
\text{elipse imaginaria}.
\]

Además:

\[
\boxed{
I_3\neq0,\ I_2<0
\Rightarrow
\text{hipérbola}
}
\]

\[
\boxed{
I_3\neq0,\ I_2=0
\Rightarrow
\text{parábola}.
}
\]

Si:

\[
I_3=0,
\]

la cónica es degenerada:

\[
\boxed{
\begin{cases}
I_2>0 & \text{punto},\\
I_2<0 & \text{dos rectas secantes},\\
I_2=0 & \text{dos rectas paralelas}.
\end{cases}
}
\]

\section{Cuádricas}

Una cuádrica es el lugar geométrico de los puntos de $\R^3$ que verifican una ecuación de segundo grado en tres variables:

\[
\boxed{
\begin{aligned}
&a_{11}x^2+a_{22}y^2+a_{33}z^2
+2a_{12}xy+2a_{13}xz+2a_{23}yz\\
&\qquad
+2a_1x+2a_2y+2a_3z+a_0=0.
\end{aligned}
}
\]

Con:

\[
X=
\begin{pmatrix}
x\\y\\z
\end{pmatrix},
\]

\[
A=
\begin{pmatrix}
a_{11}&a_{12}&a_{13}\\
a_{12}&a_{22}&a_{23}\\
a_{13}&a_{23}&a_{33}
\end{pmatrix},
\]

\[
B=
\begin{pmatrix}
2a_1&2a_2&2a_3
\end{pmatrix},
\]

se escribe:

\[
\boxed{
X^tAX+BX+a_0=0.
}
\]

En coordenadas homogéneas:

\[
\widetilde A=
\begin{pmatrix}
a_0&a_1&a_2&a_3\\
a_1&a_{11}&a_{12}&a_{13}\\
a_2&a_{12}&a_{22}&a_{23}\\
a_3&a_{13}&a_{23}&a_{33}
\end{pmatrix}.
\]

\section{Ecuación reducida de una cuádrica}

El procedimiento es análogo al de las cónicas.

\subsection{Rotación}

Se diagonaliza ortogonalmente $A$:

\[
\boxed{
D=P^tAP
=
\diag(\lambda_1,\lambda_2,\lambda_3).
}
\]

Con:

\[
X=PX',
\]

la ecuación pasa a:

\[
\lambda_1x'^2+\lambda_2y'^2+\lambda_3z'^2
+2b_1x'+2b_2y'+2b_3z'+a_0=0.
\]

\subsection{Traslación}

Completando cuadrados se eliminan los términos lineales asociados a autovalores no nulos.

Así se obtiene la ecuación reducida.

\section{Clasificación de cuádricas}

\subsection{Tres autovalores no nulos}

Si la ecuación reducida es:

\[
\lambda_1x^2+\lambda_2y^2+\lambda_3z^2+c=0,
\]

pueden aparecer:

\begin{itemize}
\item elipsoide real;
\item elipsoide imaginario;
\item hiperboloide de una hoja;
\item hiperboloide de dos hojas;
\item cono real;
\item cono imaginario.
\end{itemize}

Formas reducidas representativas:

\[
\boxed{
\frac{x^2}{a^2}
+
\frac{y^2}{b^2}
+
\frac{z^2}{c^2}
=
1
}
\]

elipsoide real,

\[
\boxed{
\frac{x^2}{a^2}
+
\frac{y^2}{b^2}
-
\frac{z^2}{c^2}
=
1
}
\]

hiperboloide de una hoja,

\[
\boxed{
\frac{x^2}{a^2}
+
\frac{y^2}{b^2}
-
\frac{z^2}{c^2}
=
-1
}
\]

hiperboloide de dos hojas,

\[
\boxed{
\frac{x^2}{a^2}
+
\frac{y^2}{b^2}
-
\frac{z^2}{c^2}
=
0
}
\]

cono real.

\subsection{Un autovalor nulo}

Si:

\[
\lambda_1=0,
\]

la ecuación tiene la forma:

\[
\lambda_2y'^2+\lambda_3z'^2+2b_1x'+b=0.
\]

Si:

\[
b_1\neq0,
\]

se obtiene:

\[
\boxed{
\begin{cases}
\lambda_2\lambda_3>0
& \text{paraboloide elíptico},\\
\lambda_2\lambda_3<0
& \text{paraboloide hiperbólico}.
\end{cases}
}
\]

Si:

\[
b_1=0,
\]

se obtienen cilindros o pares de planos.

\subsection{Dos autovalores nulos}

La ecuación queda:

\[
\lambda_3z'^2+2b_1x'+2b_2y'+b=0.
\]

Según los coeficientes, puede representar:

\begin{itemize}
\item un cilindro parabólico;
\item dos planos paralelos;
\item dos planos paralelos imaginarios;
\item un plano doble.
\end{itemize}

\section{Invariantes métricos de una cuádrica}

Se definen:

\[
\boxed{
I_4=\det\widetilde A
}
\]

\[
\boxed{
I_3=\det A
}
\]

\[
\boxed{
I_2=
\begin{vmatrix}
a_{11}&a_{12}\\
a_{12}&a_{22}
\end{vmatrix}
+
\begin{vmatrix}
a_{11}&a_{13}\\
a_{13}&a_{33}
\end{vmatrix}
+
\begin{vmatrix}
a_{22}&a_{23}\\
a_{23}&a_{33}
\end{vmatrix}
}
\]

\[
\boxed{
I_1=\tr A.
}
\]

La clasificación por invariantes de la presentación distingue:

\begin{itemize}
\item elipsoides reales e imaginarios;
\item hiperboloides de una y dos hojas;
\item paraboloides elípticos e hiperbólicos;
\item conos;
\item cilindros;
\item pares de planos.
\end{itemize}

En particular:

\[
I_4\neq0,\quad I_3=0
\]

corresponde a un paraboloide, y el signo de $I_4$ junto con los demás invariantes permite distinguir el tipo.

Si:

\[
I_4=0,
\]

la cuádrica es degenerada.

\end{document}
